\documentclass{article}

% Language setting
% Replace `english' with e.g. `spanish' to change the document language
\usepackage[czech]{babel}
\usepackage{amsthm,thmtools,xcolor,amsmath,amssymb, mathtools}

% Set page size and margins
% Replace `letterpaper' with `a4paper' for UK/EU standard size
\usepackage[a4paper,top=2cm,bottom=2cm,left=3cm,right=3cm,marginparwidth=1.75cm]{geometry}

% Useful packages
\usepackage{amsmath}
\usepackage{enumerate}
\usepackage{graphicx}
\usepackage[colorlinks=true, allcolors=blue]{hyperref}
\graphicspath{images/}

%\usepackage{tikzit}
%\input{default.tikzstyles}

\title{DÚ Lineární algebra -- Sada 1}
\author{Jan Romanovský}

\begin{document}
\maketitle

\textbf{(1.1)} Pravou stranu výrazu pomocí vlastností skalárního součinu a komplexního sdružení upravíme:

$$\frac{1}{4}(||\vec{u}+\vec{v}||^2-||\vec{u}-\vec v||^2) = \frac{1}{4}(\langle\vec{u}+\vec{v}, \vec{u}+\vec{v}\rangle-\langle \vec{u}-\vec{v}, \vec{u}-\vec{v}\rangle)= \frac{1}{4}(\langle \vec{u}, \vec{u} \rangle+\langle \vec{u}, \vec{v} \rangle+ \langle \vec{v}, \vec{u} \rangle+\langle \vec{v}, \vec{v} \rangle - \langle \vec{u}, \vec{u} \rangle -$$
$$ \langle \vec{u}, -\vec{v} \rangle-\langle -\vec{v}, \vec{u} \rangle-\langle -\vec{v}, -\vec{v} \rangle)= \frac{1}{4}(2\langle \vec{u}, \vec{v}\rangle+2\langle \vec{v}, \vec{u} \rangle ) = \frac{1}{4}(2\langle \vec{u}, \vec{v} \rangle+\overline{\langle \vec{u}, \vec{v} \rangle})=\frac{1}{4}(2\Re \langle \vec{u}, \vec{v} \rangle+2 \Im \langle \vec{u}, \vec{v} \rangle+$$
$$+2 \Re \langle \vec{u}, \vec{v} \rangle -2\Im \langle \vec{u}, \vec{v} \rangle)= \Re \langle \vec{u}, \vec{v} \rangle$$

A dostali jsme přesně výraz na levé straně.
\end{document}
